% ==========================================
% BAB IV RENCANA EVALUASI
% ==========================================
\cleardoublepage
\chapter{RENCANA EVALUASI}
\label{chap:rencana-evaluasi}

Evaluasi dalam penelitian ini bertujuan untuk menilai kualitas, kelengkapan, dan kelayakan artefak arsitektur transformasi digital Kemhan RI yang telah dirancang.
Evaluasi difokuskan pada penilaian terhadap produk arsitektur (\textit{architecture products}) sebelum implementasi.

\section{Metode Evaluasi}
Evaluasi dilakukan melalui dua pendekatan yang saling melengkapi, yaitu:
\begin{enumerate}
  \item \textit{Enterprise Architecture Scorecard}\texttrademark{}.
  \item \textit{Expert Judgement} melalui \textit{Focus Group Discussion} (FGD).
\end{enumerate}
Metode evaluasi pertama memberikan penilaian kuantitatif terstruktur yang sejalan dengan evaluasi teoritis, sedangkan metode evaluasi kedua memberikan validasi kualitatif yang sejalan dengan evaluasi praktis.

\section{Evaluasi Menggunakan \textit{Enterprise Architecture Scorecard}\texttrademark}
\textit{Enterprise Architecture Scorecard}\texttrademark{} merupakan instrumen penilaian arsitektur \textit{enterprise} yang dikembangkan oleh Jaap Schekkerman (\textit{Institute For Enterprise Architecture Developments}/IFEAD) untuk menilai kualitas suatu artefak \textit{Enterprise Architecture} berdasarkan seperangkat kriteria yang mencerminkan karakteristik arsitektur yang baik.

Tujuan:
\begin{itemize}
  \item Mengukur tingkat kelengkapan dan kualitas rancangan arsitektur secara terstruktur dan terukur dalam bentuk nilai persentase.
  \item Mengidentifikasi aspek arsitektur yang telah lengkap maupun yang masih memerlukan penyempurnaan.
\end{itemize}
Pada penelitian ini EA Scorecard\texttrademark{} digunakan untuk mengevaluasi \textit{Framework} TOGAF Pertahanan Indonesia beserta artefak arsitektur yang dihasilkan.

\subsection{Dimensi Penilaian EA Scorecard\texttrademark}
Penilaian menggunakan \textit{Enterprise Architecture Scorecard}\texttrademark{} dilakukan berdasarkan sejumlah dimensi yang merepresentasikan karakteristik utama dari arsitektur \textit{enterprise} yang baik.
Dimensi-dimensi tersebut dipilih untuk mengevaluasi sejauh mana rancangan arsitektur transformasi digital Kemhan RI mampu memenuhi kebutuhan organisasi, mendukung tujuan strategis pertahanan, serta mengakomodasi prinsip-prinsip \textit{smart defense} yang telah dirumuskan pada tahap analisis dan desain.

\begin{longtable}{|c|p{3.6cm}|p{8cm}|}
\caption{Dimensi Penilaian EA Scorecard\texttrademark}
\label{tab:dimensi-scorecard} \\
\hline
\textbf{No.} & \textbf{Dimensi} & \textbf{Deskripsi} \\
\hline
\endfirsthead
\hline
\textbf{No.} & \textbf{Dimensi} & \textbf{Deskripsi} \\
\hline
\endhead
1 & \textit{Business Alignment} & Kesesuaian arsitektur dengan visi, misi, dan strategi Kemhan. \\
\hline
2 & \textit{Completeness} & Kelengkapan artefak arsitektur yang dihasilkan. \\
\hline
3 & \textit{Integration} & Kemampuan mendukung integrasi dan interoperabilitas. \\
\hline
4 & \textit{Data Management} & Dukungan terhadap tata kelola dan integrasi data. \\
\hline
5 & \textit{Technology Readiness} & Kesesuaian teknologi dengan kebutuhan organisasi. \\
\hline
6 & \textit{Security Architecture} & Kemampuan mendukung keamanan siber dan perlindungan data. \\
\hline
7 & \textit{Governance} & Dukungan terhadap tata kelola dan kepatuhan regulasi. \\
\hline
8 & \textit{Scalability \& Flexibility} & Kemampuan beradaptasi terhadap perubahan teknologi dan ancaman. \\
\hline
9 & \textit{Smart Defense Alignment} & Kesesuaian dengan kebutuhan \textit{Smart Defense} dan C5ISR. \\
\hline
\end{longtable}

\subsection{Skala Penilaian}
Penilaian dilakukan menggunakan skala Likert 1--5 untuk mengukur tingkat kesesuaian setiap dimensi evaluasi terhadap rancangan arsitektur yang dihasilkan.
Skala ini memungkinkan evaluator memberikan penilaian secara objektif terhadap kualitas artefak arsitektur berdasarkan indikator yang telah ditetapkan.
Hasil penilaian dari seluruh evaluator kemudian diolah untuk memperoleh nilai rata-rata yang digunakan sebagai dasar dalam menentukan tingkat kualitas dan kelayakan arsitektur transformasi digital yang dirancang.

\begin{longtable}{|c|p{6cm}|}
\caption{Skala Penilaian Scorecard\texttrademark}
\label{tab:skala-penilaian} \\
\hline
\textbf{Nilai} & \textbf{Kriteria} \\
\hline
\endfirsthead
\hline
\textbf{Nilai} & \textbf{Kriteria} \\
\hline
\endhead
1 & Sangat Tidak Sesuai \\
\hline
2 & Tidak Sesuai \\
\hline
3 & Cukup Sesuai \\
\hline
4 & Sesuai \\
\hline
5 & Sangat Sesuai \\
\hline
\end{longtable}

\subsection{Interpretasi Hasil}
Hasil evaluasi \textit{Enterprise Architecture Scorecard}\texttrademark{} diinterpretasikan berdasarkan nilai rata-rata yang diperoleh dari seluruh dimensi penilaian.
Nilai tersebut menggambarkan tingkat kualitas dan kelayakan artefak arsitektur transformasi digital pertahanan yang dihasilkan.
Semakin tinggi nilai rata-rata yang diperoleh, semakin baik tingkat kesesuaian arsitektur terhadap kebutuhan organisasi, prinsip \textit{Enterprise Architecture}, serta tujuan transformasi digital pertahanan.

Dalam penelitian ini, artefak arsitektur dinyatakan layak apabila memperoleh nilai rata-rata minimal pada kategori \textbf{Baik}, yang menunjukkan bahwa rancangan arsitektur telah memenuhi sebagian besar kriteria \textit{Enterprise Architecture} dan dapat direkomendasikan untuk tahap implementasi lebih lanjut.

\begin{longtable}{|c|p{6cm}|}
\caption{Interpretasi Hasil}
\label{tab:interpretasi-hasil} \\
\hline
\textbf{Rentang Nilai} & \textbf{Interpretasi} \\
\hline
\endfirsthead
\hline
\textbf{Rentang Nilai} & \textbf{Interpretasi} \\
\hline
\endhead
4,20 -- 5,00 & Sangat Baik \\
\hline
3,40 -- 4,19 & Baik \\
\hline
2,60 -- 3,39 & Cukup \\
\hline
1,80 -- 2,59 & Kurang \\
\hline
1,00 -- 1,79 & Sangat Kurang \\
\hline
\end{longtable}

\section{Evaluasi oleh \textit{Expert Judgement}}
Evaluasi ini dilakukan dengan tujuan untuk memperoleh penilaian dari para ahli dan pemangku kepentingan yang berada pada sektor pertahanan terkait kelayakan, relevansi, dan implementabilitas desain arsitektur transformasi digital yang dihasilkan.

Peserta:
\begin{itemize}
  \item Pakar akademik bidang \textit{enterprise architecture} dan sistem informasi.
  \item Praktisi dan pejabat Kemhan.
\end{itemize}

Untuk melengkapi hasil evaluasi menggunakan \textit{Enterprise Architecture Scorecard}\texttrademark{}, penelitian ini juga menerapkan metode \textit{Expert Judgement} yang melibatkan akademisi dan praktisi yang memiliki kompetensi di bidang transformasi digital, \textit{enterprise architecture}, dan sektor pertahanan.
Evaluasi ini bertujuan untuk memperoleh masukan, validasi, dan penilaian profesional terhadap rancangan TOGAF Pertahanan Indonesia yang telah dikembangkan.
Aspek penilaian disusun berdasarkan komponen utama arsitektur yang dirancang, meliputi aspek strategis, organisasi, data dan sistem informasi, teknologi, serta C5ISR.
Setiap aspek dievaluasi berdasarkan kriteria yang merepresentasikan tingkat kesesuaian arsitektur terhadap kebutuhan transformasi digital Kemhan RI, konsep \textit{smart defense}, regulasi yang berlaku, serta kelayakan implementasinya di lingkungan pertahanan.
Adapun aspek dan kriteria penilaian yang digunakan dalam proses \textit{Expert Judgement} disajikan pada Tabel \ref{tab:aspek-penilaian}.

\begin{longtable}{|c|p{3.2cm}|p{7cm}|c|}
\caption{Aspek Penilaian \textit{Expert Judgement}}
\label{tab:aspek-penilaian} \\
\hline
\textbf{No.} & \textbf{Dimensi} & \textbf{Kriteria} & \textbf{Nilai} \\
\hline
\endfirsthead
\hline
\textbf{No.} & \textbf{Dimensi} & \textbf{Kriteria} & \textbf{Nilai} \\
\hline
\endhead
1 & Aspek Strategis & Kesesuaian dengan visi transformasi digital Kemhan. & \\
\hline
2 & & Kesesuaian dengan \textit{Smart Defense}. & \\
\hline
3 & & Kesesuaian dengan regulasi nasional. & \\
\hline
4 & Aspek Organisasi & Kesesuaian struktur organisasi target. & \\
\hline
5 & & Kelayakan pembentukan DTO (\textit{Digital Transformation Office}). & \\
\hline
6 & & Kesiapan tata kelola transformasi digital. & \\
\hline
7 & Aspek Data dan Sistem & Integrasi data pertahanan. & \\
\hline
8 & & Interoperabilitas aplikasi. & \\
\hline
9 & & Tata kelola data pertahanan. & \\
\hline
10 & Aspek Teknologi & Kelayakan teknologi yang diusulkan. & \\
\hline
11 & & Kesiapan implementasi. & \\
\hline
12 & & Keamanan siber. & \\
\hline
13 & Aspek C5ISR & Dukungan terhadap \textit{Command and Control}. & \\
\hline
14 & & Dukungan terhadap interoperabilitas lintas matra. & \\
\hline
15 & & Dukungan terhadap \textit{Network Centric Warfare}. & \\
\hline
\end{longtable}

Mekanisme pelaksanaan:
\begin{itemize}
  \item Presentasi hasil desain arsitektur.
  \item Sesi diskusi dan tanya jawab.
  \item Pengisian instrumen penilaian.
  \item Pengumpulan rekomendasi dan perbaikan.
  \item Penyusunan hasil evaluasi.
\end{itemize}

\section{Kriteria Keberhasilan Evaluasi}
Artefak arsitektur dinyatakan valid apabila:
\begin{itemize}
  \item Nilai EA Scorecard\texttrademark{} memperoleh rata-rata $\geq$ 4,00.
  \item Minimal 80\% peserta FGD menyatakan desain arsitektur layak diterapkan.
  \item Tidak terdapat temuan kritis yang menghambat implementasi.
  \item Seluruh rekomendasi perbaikan dapat diakomodasi dalam penyempurnaan desain.
\end{itemize}
